\documentclass[11pt]{article}

% ---------------- Packages ----------------
\usepackage[a4paper,margin=1in]{geometry}
\usepackage{amsmath,amssymb,amsfonts}
\usepackage{mathtools}
\usepackage{bm}
\usepackage{booktabs}
\usepackage{array}
\usepackage{longtable}
\usepackage{tabularx}
\usepackage{enumitem}
\usepackage{graphicx}
\usepackage{xcolor}
\usepackage[colorlinks=true,linkcolor=blue,citecolor=blue,urlcolor=blue]{hyperref}
\usepackage[activate={true,nocompatibility},final,tracking=true,kerning=true,spacing=true,factor=1100,stretch=10,shrink=10,expansion=false,protrusion=true]{microtype}
\usepackage{setspace}
\usepackage{titlesec}
\usepackage{fancyhdr}
\usepackage{footmisc}
\usepackage{parskip}
\usepackage{url}
\usepackage{xspace}
\usepackage{authblk}
\usepackage[numbers,sort&compress]{natbib}

\titleformat{\section}{\Large\bfseries}{\thesection}{0.8em}{}
\titleformat{\subsection}{\large\bfseries}{\thesubsection}{0.8em}{}
\titleformat{\subsubsection}{\normalsize\bfseries\itshape}{\thesubsubsection}{0.6em}{}

\pagestyle{fancy}
\fancyhf{}
\fancyhead[L]{\small Mouse kidney spaceflight network contrasts}
\fancyhead[R]{\small \thepage}
\renewcommand{\headrulewidth}{0.4pt}

\newcommand{\code}[1]{\texttt{#1}}
\newcommand{\pkg}[1]{\textsf{#1}}
\newcommand{\gene}[1]{\textit{#1}}
\newcommand{\RRRM}{RRRM-2\xspace}

\title{\textbf{ISS-T-Aligned Extracellular-Matrix and Tubular-Transport Remodeling in Mouse Kidney Spaceflight Transcriptomes}\\[0.35em]
\large A multi-resolution network-contrast analysis across RRRM-2 and external OSDR cohorts}
\author[1]{Ibrahim Shahid}
\author[1,*]{}
\affil[1]{Independent research, manuscript in preparation}
\affil[*]{Senior author placeholder pending PI involvement}
\date{May 2026}

\begin{document}
\maketitle

% =====================================================================
\begin{abstract}
\noindent
Spaceflight alters renal physiology and is associated with kidney-stone risk, tubulopathy, and extracellular-matrix remodeling. NASA Rodent Research Reference Mission 2 (RRRM-2; OSD-771) provides a rare factorial mouse kidney transcriptome design with Age, mission Arm, and Environment Group represented in the same experiment, enabling questions about whether spaceflight modifies age-associated kidney transcriptional programs differently in the ISS terminal dissection arm (ISS-T) and the live animal return arm (LAR). We analyzed residualized RRRM-2 kidney expression with a multi-resolution network-contrast framework that combined WGCNA module eigengenes, curated pathway scores, LIONESS-derived sample-specific node summaries, arm-specific flight vectors, contrast-level recurrence against external OSDR cohorts, and projection onto an independently estimated Tabula Muris Senis kidney aging axis. The central within-cohort age-rewiring contrast was treated explicitly as a vector problem: for each arm, control aging was estimated as Old GC minus Young GC and compared to the corresponding flight aging vector. Because RRRM-2 contains only $n=5$ mice per Age $\times$ Arm $\times$ Environment Group cell, we first quantified bootstrap angular stability of the control-aging vector before interpreting any age-rewiring decomposition.

The production run (\code{run\_20260518\_201823\_2500g}) does not support a primary claim that RRRM-2 alone resolves age-dependent network rewiring at module or pathway scale. Primary module/pathway control-aging directions were not consistently stable across arms: ISS-T pathway lower 95\% cosine bound was 0.248, ISS-T module was 0.056, and LAR pathway was 0.345, all below the pre-specified 0.40 lower-bound criterion, although secondary gene and LIONESS-node summaries were stable in both arms. The positive result emerges instead at the cross-cohort flight-vector level. The RRRM-2 ISS-T pathway flight vector aligned with an independent OSD-513 flight vector (male C57BL/6J kidney; 38-day flight; $n=9$ FLT and $n=9$ HGC) with cosine 0.641 and bootstrap 95\% CI 0.373 to 0.811. The RRRM-2 LAR vector was anti-aligned with OSD-513 (cosine $-0.511$, CI $-0.712$ to $-0.216$). Pathway decomposition showed that the ISS-T recurrence was driven primarily by concordant up-shifts in ECM remodeling and fibrosis and concordant down-shifts in DCT/NCC-WNK transport biology. OSD-253 C57BL/6J comparisons showed qualitatively similar ECM/fibrosis-positive direction under both original and white-LED rerun ground-control definitions, but confidence intervals crossed zero and the cohort was retained as context sensitivity because of light and read-length control-definition caveats. Projection onto the Tabula Muris Senis kidney aging axis did not support a strong aging-mimicry interpretation: ISS-T was essentially orthogonal to the external aging axis ($\beta=-0.0027$, cosine $=-0.0187$), while LAR showed only a small positive projection ($\beta=0.0162$, cosine $=0.175$).

Together, these results support a constrained biological claim. RRRM-2 does not by itself establish broad age-dependent kidney network rewiring after spaceflight. Instead, it identifies an ISS-T-specific flight direction that recurs in an independent OSDR kidney cohort and is dominated by ECM/fibrotic remodeling and tubular-transport changes. WGCNA module results, including the previously identified grey60 ECM/cell-migration module, are best interpreted as structural annotation of this remodeling biology rather than as the primary inferential basis. The study illustrates how small-cohort spaceflight transcriptomes can still yield biological signal when factorial contrasts, sample-specific networks, module analysis, and external contrast-level recurrence are integrated with explicit attention to what each method can and cannot identify.
\end{abstract}

\vspace{0.3em}
\noindent\textbf{Keywords:} spaceflight transcriptomics; mouse kidney; OSD-771; OSD-513; OSD-253; RRRM-2; ISS-T; LAR; LIONESS; WGCNA; Tabula Muris Senis; extracellular matrix; fibrosis; tubular transport.

\clearpage
\tableofcontents
\clearpage

% =====================================================================
\section{Introduction}

\subsection{Spaceflight kidney risk and the need for transcriptome-scale renal biology}

Renal risk has been recognized as a practical challenge for human spaceflight for more than half a century. Early work predicted that prolonged weightlessness could increase urolithiasis risk through skeletal calcium loss \citep{cockett1962astronautical}, and later astronaut studies confirmed that kidney-stone risk remains elevated after long-duration missions \citep{pietrzyk2007renal,smith2014men,kassemi2021numerical}. Recent multi-omic and physiological work has broadened the problem from downstream stone risk to primary kidney remodeling, including tubulopathy, altered transport physiology, mitochondrial stress, lipid dysregulation, and extracellular-matrix (ECM) responses \citep{daSilveira2020comprehensive,siew2024cosmic,chouker2025impact}. The kidney is therefore not simply a passive target of systemic spaceflight effects; it is an active tissue in which flight-associated remodeling may contribute to clinically relevant renal dysfunction.

Mouse spaceflight datasets in the NASA Open Science Data Repository (OSDR) are central to this question. They provide controlled genetic backgrounds, flight and ground-control groups, multiple mission durations, and tissue-level transcriptomes. However, most rodent spaceflight omics studies remain limited by small per-condition sample sizes, mission-specific technical conditions, and heterogeneous metadata. These constraints do not invalidate the datasets, but they require analyses that distinguish what can be estimated within a single cohort from what must be evaluated across independent cohorts.

\subsection{RRRM-2 enables an unusually specific question}

NASA Rodent Research Reference Mission 2 (RRRM-2; OSD-771) is valuable because it contains a full factorial design in female C57BL/6NTac kidney. Age (Young, Old), mission Arm (ISS-T, LAR), and Environment Group (FLT, Habitat Ground Control, Vivarium Control, Basal) are represented with $n=5$ biological replicates per cell. This structure makes RRRM-2 better suited than many OSDR cohorts for asking whether the kidney response to spaceflight depends on age and arm.

The motivating question is not merely whether a gene or module is age-associated, flight-associated, or enriched for a pathway. It is whether spaceflight changes the direction of age-associated kidney transcriptional programs. Within each arm $a \in \{\mathrm{ISS\mbox{-}T},\mathrm{LAR}\}$, the relevant control-aging and flight-aging vectors are
\begin{equation}
A^{a}_{GC} = \bar{X}_{\mathrm{Old},a,GC} - \bar{X}_{\mathrm{Young},a,GC},
\qquad
A^{a}_{FLT} = \bar{X}_{\mathrm{Old},a,FLT} - \bar{X}_{\mathrm{Young},a,FLT}.
\end{equation}
The comparison between $A^{a}_{FLT}$ and $A^{a}_{GC}$ asks whether flight amplifies, dampens, reverses, or redirects the age-associated direction. This is a factorial contrast problem. It is related to, but not identical with, differential expression, WGCNA module-trait association, or differential correlation.

\subsection{Why multiple network representations are needed}

No single network method directly resolves all layers of this design. WGCNA identifies co-expression modules and module eigengene shifts \citep{langfelder2008wgcna}; it is useful for finding stable biological programs but does not naturally test edge-wise or node-wise factorial interactions. LIONESS estimates sample-specific network contributions \citep{kuijjer2019estimating}; it provides per-mouse network features, but in small cohorts its predictive value must be benchmarked against expression-only baselines and interpreted cautiously. Direct differential correlation tests group-level topology changes, but separate $n=5$ networks are noisy. Pathway-level vectors are lower-dimensional and interpretable, but they compress gene-level heterogeneity.

We therefore combined methods rather than treating one as definitive. WGCNA modules were used as a module-scale structural layer. Curated pathway scores provided interpretable low-dimensional flight and aging vectors. LIONESS-derived node summaries supplied a sample-specific network layer. External OSDR cohorts were used only at contrast level, avoiding raw expression pooling across missions. Finally, an independent Tabula Muris Senis kidney aging axis was used to test whether flight effects align with kidney aging without relying entirely on RRRM-2's small within-arm control-aging estimates.

\subsection{Study objective}

The objective of this study is to determine what RRRM-2 can support about age- and arm-dependent kidney response to spaceflight when analyzed jointly with external OSDR cohorts. We ask three questions. First, are within-RRRM control-aging directions stable enough to interpret age-rewiring contrasts at module/pathway scale? Second, do RRRM-2 arm-specific flight directions recur in independent OSDR kidney cohorts? Third, do the observed flight directions resemble an independent kidney aging axis, or are they better understood as remodeling responses distinct from aging mimicry?

% =====================================================================
\section{Methods}

\subsection{RRRM-2 discovery cohort}

RRRM-2 (OSD-771) was analyzed as the discovery and factorial-bridge cohort. The design contains female C57BL/6NTac mouse kidney samples across Age (Young and Old), Arm (ISS-T and LAR), and Environment Group (FLT, Habitat Ground Control, Vivarium Control, and Basal), with $n=5$ per Age $\times$ Arm $\times$ Environment Group cell. For the primary contrast-vector analyses, Habitat Ground Control labels were normalized to GC and compared with FLT. Vivarium and Basal samples were retained in preprocessing resources but were not used in the FLT-versus-GC vectors reported here.

ISS-T and LAR were treated as arm labels rather than pure causal mechanisms. ISS-T mice were sacrificed on orbit, whereas LAR mice returned alive and were sacrificed after recovery on Earth. Consequently, arm differences may reflect flight state, recovery window, and tissue collection protocol. This is a known issue in rodent spaceflight transcriptomics, where sample collection and preservation can introduce large transcriptomic effects \citep{galazka2020rnaseq}.

\subsection{Expression preprocessing and residualized feature matrix}

The analysis used the repository's residualized expression matrix \code{Rtech.tsv.gz} and matched metadata \code{meta\_phase1.tsv.gz}. Upstream preprocessing followed the project pipeline: count filtering, variance-stabilizing transformation, cell-type deconvolution against a Tabula Muris Senis kidney reference, and residualization against technical and cellular-composition covariates. The residual matrix was indexed by genes and samples and served as the basis for pathway scores, WGCNA eigengene analysis, and LIONESS feature integration.

\subsection{Network and pathway feature resolutions}

Four feature resolutions were evaluated.

\textbf{Module resolution.} WGCNA module eigengenes from the existing WGCNA run were used as module-level features. WGCNA modules capture unsupervised co-expression programs and provide a natural bridge to biological annotation.

\textbf{Pathway resolution.} Curated kidney-relevant gene sets from \code{config/gene\_sets.yaml} were converted to sample-level pathway scores by averaging residualized expression across mapped genes. The production pathway panel contained nine features: DCT/NCC-WNK, ECM remodeling, oxidative stress, calcium handling, lipid metabolism, inflammation, apoptosis, ion transport, and fibrosis.

\textbf{Gene and LIONESS-node resolution.} Sample-specific LIONESS edge matrices were loaded from the Phase 2 network output and aligned to the edge index and sample metadata. Absolute incident sample-specific edge strengths were averaged per gene to produce node-level features. The \code{gene} and \code{lioness\_node} labels therefore both represent LIONESS-derived node summaries in the production run, with the latter retained as an explicit sample-specific-network resolution.

\subsection{Within-cohort age-vector stability and decomposition}

For each arm and resolution, the control-aging vector $A_{GC}$ was estimated as Old GC minus Young GC. Bootstrap resampling was stratified by Age, Arm, and Environment Group. Each bootstrap replicate rebuilt $A_{GC}$ and computed cosine similarity to the full-sample vector. Module and pathway resolutions were treated as primary because they are the intended scales for interpretable biology claims; gene and LIONESS-node resolutions were treated as secondary.

If primary control-aging vectors were sufficiently stable, the planned within-RRRM decomposition compared $A_{FLT}$ with $A_{GC}$ using projection coefficient $\beta$, cosine similarity, and residual fraction. The projection coefficient
\begin{equation}
\beta = \frac{A_{FLT}\cdot A_{GC}}{\lVert A_{GC}\rVert^2}
\end{equation}
quantifies amplification ($\beta>1$), dampening ($0<\beta<1$), or reversal ($\beta<0$) along the control-aging axis. The residual component quantifies redirection away from that axis. In the production run, primary module/pathway stability did not pass consistently, so within-RRRM age-vector decomposition was not reported as a primary result.

\subsection{External OSDR contrast-level recurrence}

External cohorts were used at contrast level rather than by raw expression pooling. For each external study $s$, a study-specific flight vector was computed as
\begin{equation}
F^{s} = \bar{X}^{s}_{FLT} - \bar{X}^{s}_{GC}.
\end{equation}
RRRM-2 arm-specific flight vectors were computed as age-pooled FLT minus GC within each arm. External and RRRM-2 vectors were compared by cosine similarity on shared pathway features. External samples were bootstrapped within FLT and GC strata to estimate a percentile interval.

OSD-513 was used as the directional-recurrence cohort because it provides a powered C57BL/6J kidney flight comparison with $n=9$ FLT and $n=9$ Habitat Ground Control samples. OSD-253 was analyzed as a context-sensitivity cohort, not as clean validation, because its original and rerun ground-control definitions trade off light-wavelength and read-length/platform compatibility. For OSD-253, C57BL/6J comparisons were computed against both original blue-light ground controls and the 2021 white-LED rerun controls.

\subsection{External kidney aging-axis projection}

An independent kidney aging axis was built from Tabula Muris Senis female kidney data \citep{tabula2020}. Cells were aggregated to donor-level pseudobulk logCPM profiles, and the primary aging axis was defined as 18-month minus 3-month donor expression. RRRM-2 age-pooled arm-specific flight vectors were projected onto this external aging axis using the same projection statistics as above. Bootstrap resampling propagated RRRM-2 sample uncertainty and, after donor-level source files were rebuilt, Tabula Muris Senis donor uncertainty.

\subsection{Secondary WGCNA and candidate-overlap analyses}

WGCNA module results were used as structural annotation. Previously identified module patterns, including grey60 ECM/cell-migration biology, were compared against LIONESS candidate genes, silent-shifter candidate distributions, and module eigengene simple effects. Fisher's exact tests evaluated candidate enrichment across modules with Benjamini-Hochberg correction \citep{benjamini1995controlling}. These analyses were used to determine whether gene-level sample-specific network candidates, silent shifters, and WGCNA module activity converged on a single module or instead represented method-specific views of the response.

\subsection{Software and reproducibility}

The production run was launched with:
\begin{center}
\code{venv/bin/python -m src.run\_all\_phases --updated-pipeline}.
\end{center}
The run ending in \code{201823\_2500g} generated the v3 result artifacts under its \code{contrast\_vectors} output directory. It used 2000 bootstrap iterations for stability and external projection. The run manifest records input hashes for the residualized expression matrix, metadata, WGCNA module directory, LIONESS network directory, and configuration file. Focused tests for contrast decomposition, stability testing, external aging-axis projection, cross-OSDR alignment, decision logic, and LIONESS methodology passed after implementation.

% =====================================================================
\section{Results}

\subsection{WGCNA, LIONESS candidates, and silent shifters do not collapse into one module-level story}

WGCNA previously identified grey60 as a biologically coherent ECM/cell-migration module with an ISS-T-specific eigengene pattern. That module remains important annotation, but follow-up cross-method tests did not support making it the single organizing result. The 11 ISS-T-Young LIONESS factorial edge-rewiring candidates did not concentrate in grey60; grey60 contained 0/11 candidates. Candidate genes instead landed across purple (3 genes: \gene{Nrip3}, \gene{Nupr1}, \gene{Aldh1a3}), turquoise (3: \gene{Sap25}, \gene{Cox6b2}, \gene{Slc2a9}), yellow (2: \gene{C1qtnf3}, \gene{Epdr1}), black (1: \gene{Ache}), brown (1: \gene{Ppp1r1b}), and one off-panel Ensembl feature. Purple was the strongest candidate enrichment (fold enrichment $10.2\times$), but its BH-adjusted value was marginal ($q=0.054$).

Silent-shifter candidates showed a similarly distributed pattern. ISS-T-Young silent shifters enriched midnightblue (11/185; $q=0.0031$), ISS-T-Old enriched purple (13/171; $q=0.030$), LAR-Young enriched grey/unassigned and purple ($q=0.047$ and $0.092$), and LAR-Old had no module below $q<0.10$. Module eigengene simple effects also showed that grey60 was not unique: grey60, royalblue, salmon, and pink all displayed significant ISS-T-specific eigengene patterns. Thus the WGCNA result supports ECM/fibrosis annotation, but it does not provide a unique module-level convergence point for every network-derived signal.

\subsection{Within-RRRM control-aging directions are unstable at primary module/pathway scale}

The first production analysis evaluated whether within-arm control-aging vectors could support the planned age-rewiring decomposition. Table~\ref{tab:stability} reports bootstrap angular stability of $A_{GC}$ by arm and resolution. ISS-T pathway, ISS-T module, and LAR pathway failed the lower-bound stability criterion, whereas LAR module passed. Secondary gene and LIONESS-node summaries passed in both arms.

\begin{table}[h]
\centering\small
\begin{tabular}{llrrrrc}
\toprule
Arm & Resolution & Full norm & Median cosine & Lower CI & Required lower & Passed \\
\midrule
ISS-T & pathway & 0.167 & 0.787 & 0.248 & 0.400 & no \\
ISS-T & module & 0.325 & 0.750 & 0.056 & 0.400 & no \\
ISS-T & gene & 10.152 & 0.746 & 0.545 & 0.200 & yes \\
ISS-T & LIONESS-node & 10.152 & 0.745 & 0.549 & 0.200 & yes \\
LAR & pathway & 0.221 & 0.890 & 0.345 & 0.400 & no \\
LAR & module & 0.353 & 0.774 & 0.493 & 0.400 & yes \\
LAR & gene & 9.687 & 0.742 & 0.563 & 0.200 & yes \\
LAR & LIONESS-node & 9.687 & 0.737 & 0.556 & 0.200 & yes \\
\bottomrule
\end{tabular}
\caption{Bootstrap angular stability of RRRM-2 control-aging vectors. Module and pathway are primary biological resolutions; gene and LIONESS-node are secondary sample-specific-network summaries.}
\label{tab:stability}
\end{table}

This result narrows the within-cohort claim. RRRM-2 does not support a primary module/pathway statement that spaceflight amplifies, dampens, or redirects age-associated kidney network structure. The relevant control-aging direction is too sensitive to resampling in three of four primary arm-resolution settings. The passed secondary gene and LIONESS-node stability results indicate that lower-level node summaries may still be useful for exploratory analysis, but they should not be promoted into the main age-rewiring conclusion.

\subsection{ISS-T, but not LAR, recurs in the independent OSD-513 flight vector}

Although the within-RRRM age-rewiring decomposition is not primary-claimable, the arm-specific flight vectors can still be compared to external flight contrasts. OSD-513 provided a powered independent kidney flight vector. RRRM-2 ISS-T aligned positively with OSD-513 (point cosine 0.641; median 0.624; 95\% CI 0.373 to 0.811). RRRM-2 LAR was anti-aligned with OSD-513 (point cosine $-0.511$; CI $-0.712$ to $-0.216$). OSD-253 comparisons were directionally informative but retained as context sensitivity because both available control definitions have design caveats (Table~\ref{tab:crossosdr}).

\begin{table}[h]
\centering\small
\resizebox{\textwidth}{!}{%
\begin{tabular}{lllrrrrcc}
\toprule
Study & Scenario & Arm & Point & Median & Lower CI & Upper CI & Boundary & Claimable \\
\midrule
OSD-513 & powered HGC & ISS-T & 0.641 & 0.624 & 0.373 & 0.811 & recurrence & yes \\
OSD-513 & powered HGC & LAR & -0.511 & -0.494 & -0.712 & -0.216 & recurrence & no \\
OSD-253 & original GC blue light & ISS-T & 0.674 & 0.605 & -0.345 & 0.836 & context & no \\
OSD-253 & original GC blue light & LAR & -0.679 & -0.612 & -0.809 & 0.384 & context & no \\
OSD-253 & rerun GC white light & ISS-T & 0.474 & 0.451 & -0.024 & 0.689 & context & no \\
OSD-253 & rerun GC white light & LAR & -0.318 & -0.314 & -0.578 & 0.201 & context & no \\
\bottomrule
\end{tabular}
}
\caption{Cross-OSDR pathway-vector alignment. OSD-513 is treated as directional recurrence. OSD-253 is treated as context sensitivity because of control-definition caveats.}
\label{tab:crossosdr}
\end{table}

The key biological implication is arm specificity. The ISS-T flight direction resembles an independent flight response; the LAR direction does not. This does not prove that the difference is caused only by acute on-orbit flight state, because LAR also differs in recovery and collection protocol. It does show that the terminal-flight arm captures a recurring kidney remodeling direction that is absent or reversed after live animal return.

\subsection{The recurring ISS-T direction is dominated by ECM/fibrosis and tubular-transport biology}

Pathway-level decomposition identified the drivers of the OSD-513 alignment (Table~\ref{tab:pathwaydrivers}). ECM remodeling and fibrosis were the largest positive contributors to the ISS-T cosine. DCT/NCC-WNK also contributed positively because both RRRM-2 ISS-T and OSD-513 moved downward on that pathway. Lipid metabolism and ion transport opposed the ISS-T alignment at this pathway resolution.

\begin{table}[h]
\centering\small
\resizebox{\textwidth}{!}{%
\begin{tabular}{lrrrrr}
\toprule
Pathway & RRRM ISS-T & RRRM LAR & OSD-513 & ISS-T cosine contribution & LAR cosine contribution \\
\midrule
ECM remodeling & 0.242 & -0.197 & 0.461 & 0.352 & -0.361 \\
Fibrosis & 0.240 & -0.201 & 0.398 & 0.302 & -0.319 \\
DCT/NCC-WNK & -0.075 & -0.028 & -0.370 & 0.088 & 0.041 \\
Apoptosis & 0.098 & -0.026 & 0.112 & 0.035 & -0.012 \\
Calcium handling & 0.000 & -0.034 & -0.130 & -0.000 & 0.018 \\
Oxidative stress & 0.020 & -0.062 & -0.004 & -0.000 & 0.001 \\
Inflammation & 0.058 & 0.014 & -0.075 & -0.014 & -0.004 \\
Lipid metabolism & -0.117 & 0.086 & 0.158 & -0.058 & 0.054 \\
Ion transport & 0.070 & -0.062 & -0.290 & -0.064 & 0.071 \\
\bottomrule
\end{tabular}
}
\caption{Pathway drivers of OSD-513 recurrence. RRRM columns are age-pooled FLT-minus-GC pathway effects by arm. OSD-513 is FLT-minus-GC. Cosine contributions are normalized dot-product contributions.}
\label{tab:pathwaydrivers}
\end{table}

This is the main biological result. The recurring ISS-T-aligned flight response is a kidney remodeling program characterized by ECM/fibrosis activation and tubular-transport change. This agrees with prior spaceflight kidney literature implicating ECM/TGF-$\beta$ signaling, renal remodeling, lipid/mitochondrial stress, and altered tubule physiology \citep{beheshti2018microrna,daSilveira2020comprehensive,finch2025strain,siew2024cosmic}.

\subsection{LAR anti-alignment is driven by ECM/fibrosis reversal}

The LAR anti-alignment with OSD-513 is driven primarily by ECM remodeling and fibrosis. OSD-513 shows positive ECM/fibrosis flight effects. RRRM-2 ISS-T shows the same direction. RRRM-2 LAR shows the opposite direction. This reversal explains why LAR cosine is negative despite some same-direction movement in DCT/NCC-WNK and broader ion-transport features.

Biologically, LAR may represent attenuation of a terminal-flight ECM response, a recovery-specific transcriptional state, or a composite of recovery and tissue collection differences. The current data do not separate those mechanisms. The supported statement is descriptive: LAR does not recapitulate the OSD-513-like flight remodeling direction and differs most strongly in ECM/fibrosis features.

\subsection{OSD-253 is context-compatible but not a clean validation cohort}

The OSD-253 C57BL/6J comparisons are useful as sensitivity analyses. Against original blue-light ground controls, ISS-T alignment is positive at the point estimate (0.674) but has a wide interval crossing zero. Against the white-LED rerun controls, ISS-T alignment remains positive but weaker (0.474) and again crosses zero. In both definitions, LAR directions are negative. At the pathway level, fibrosis and ECM remodeling are the strongest positive contributors to ISS-T alignment under both control definitions.

These results are compatible with the OSD-513 recurrence pattern, but they should not be described as validation. Original controls are read-length compatible but light-confounded; rerun controls are light-compatible but read-length/platform-shifted. OSD-253 is therefore best used to evaluate context dependence and to motivate a future strain-stratified analysis involving C3H/HeJ.

\subsection{External kidney aging-axis projection does not support aging mimicry}

Projection of RRRM-2 arm-specific flight vectors onto the Tabula Muris Senis kidney aging axis did not support a strong aging-mimicry interpretation (Table~\ref{tab:tms}). ISS-T was essentially orthogonal to the external aging axis. LAR had a positive confidence interval for $\beta$ and cosine, but the effect was small and most of the vector remained residual to the aging direction.

\begin{table}[h]
\centering\small
\begin{tabular}{lrrrrrr}
\toprule
Arm & Statistic & Point & Median & Lower CI & Upper CI & Interpretation \\
\midrule
ISS-T & $\beta$ & -0.0027 & -0.0023 & -0.0166 & 0.0180 & redirected \\
ISS-T & cosine & -0.0187 & -0.0178 & -0.1076 & 0.1064 & redirected \\
ISS-T & $\rho$ & 0.9998 & 0.9991 & 0.9928 & 1.0000 & redirected \\
LAR & $\beta$ & 0.0162 & 0.0110 & 0.0024 & 0.0341 & redirected \\
LAR & cosine & 0.1753 & 0.1233 & 0.0328 & 0.2296 & redirected \\
LAR & $\rho$ & 0.9845 & 0.9924 & 0.9733 & 0.9995 & redirected \\
\bottomrule
\end{tabular}
\caption{Projection of RRRM-2 age-pooled flight vectors onto an independent Tabula Muris Senis kidney female aging axis (18 months minus 3 months).}
\label{tab:tms}
\end{table}

This result is important because it prevents overstatement. The positive cross-OSDR signal should not be described as flight accelerating kidney aging. It is better described as a flight-associated remodeling and transport response that is largely distinct from the external aging axis used here.

% =====================================================================
\section{Discussion}

\subsection{Principal finding}

The central biological finding is that RRRM-2 ISS-T captures a recurring spaceflight kidney remodeling direction that is visible in an independent OSDR cohort. The direction is dominated by ECM remodeling, fibrosis, and tubular-transport changes. RRRM-2 LAR does not share this direction and is anti-aligned with OSD-513, primarily because its ECM/fibrosis effects have the opposite sign.

The central statistical boundary is equally clear: RRRM-2 alone does not support a primary module/pathway claim about age-dependent network rewiring. The age-rewiring contrast requires a stable within-arm control-aging direction, and at $n=5$ per cell that direction is not consistently stable at the primary biological resolutions. This finding does not make the dataset useless; it clarifies which claim is supported. The dataset supports arm-specific flight-vector recurrence, not standalone within-cohort aging-rewiring inference.

\subsection{Biological interpretation}

The recurring ISS-T response fits a kidney-remodeling model rather than a broad topology-rewiring model. ECM and fibrosis pathway increases are the dominant positive contributors to cross-OSDR recurrence. DCT/NCC-WNK suppression contributes in the same direction because it decreases in both RRRM-2 ISS-T and OSD-513. This combination is biologically plausible: matrix remodeling and altered tubular transport are both prominent themes in spaceflight kidney biology and in renal-stress responses more generally \citep{finch2025strain,siew2024cosmic}.

The LAR result is not simply absence of signal. LAR is anti-aligned with OSD-513 because the ECM/fibrosis component reverses. This could reflect recovery after return, a distinct post-flight adaptation, or sample-collection differences. The present design cannot distinguish these explanations, and the manuscript should avoid causal language stronger than the data allow. ``ISS-T-aligned remodeling'' is supported; ``acute on-orbit ECM activation that reverses with recovery'' remains a hypothesis.

\subsection{Relationship to WGCNA and sample-specific network analyses}

The WGCNA module layer remains useful but should not be the sole foundation of the paper. Grey60 provides a coherent ECM/cell-migration module annotation that is consistent with the pathway-level cross-OSDR result. However, LIONESS candidate genes and silent-shifter distributions do not converge uniquely on grey60, and multiple modules share ISS-T-specific eigengene patterns. The combined interpretation is therefore not ``one module explains everything.'' It is that several analytical views point toward ISS-T-associated remodeling biology, with ECM/fibrosis emerging most clearly at the pathway and WGCNA-annotation levels.

The LIONESS-derived node summaries are also informative. Their stability at secondary resolution suggests that sample-specific network-derived features may support exploratory gene/node prioritization. However, they should not override the primary module/pathway instability for the original age-rewiring question. A sensible next step is to report exploratory secondary decomposition for gene and LIONESS-node features and ask which genes drive the ISS-T-aligned ECM/fibrosis and transport direction.

\subsection{What the study does and does not claim}

This study supports the following claims. First, RRRM-2 ISS-T, but not LAR, aligns with an independent OSD-513 kidney flight vector. Second, the shared ISS-T/OSD-513 direction is dominated by ECM/fibrosis activation and tubular-transport remodeling. Third, OSD-253 C57BL/6J comparisons are qualitatively compatible with this direction but are not clean validation because of control-definition caveats. Fourth, projection onto an external kidney aging axis does not support a strong aging-mimicry interpretation. Fifth, WGCNA modules provide useful structural annotation but do not reduce the entire response to a single grey60-centered convergence story.

The study does not claim that RRRM-2 proves age-dependent network rewiring at module/pathway scale. It does not claim broad co-expression topology rewiring. It does not claim that spaceflight simply accelerates kidney aging. It does not claim that OSD-253 validates the result. It does not claim a TLR4/C3H-HeJ mechanism from the current production run, because the cross-OSDR recurrence analysis reported here used C57BL/6J OSD-253 comparisons only.

\subsection{Limitations}

Several limitations shape the interpretation. First, RRRM-2 has only $n=5$ samples per factorial cell. That design is rich but underpowered for high-dimensional contrast-of-contrasts at module/pathway scale. Second, ISS-T and LAR differ by recovery and tissue collection protocol as well as biological flight state \citep{galazka2020rnaseq}. Third, OSD-513 differs from RRRM-2 in sex, strain substrain, mission context, and other cohort features; it is an external recurrence cohort, not a perfect replication. Fourth, OSD-253 contains unavoidable control-definition tradeoffs. Fifth, the Tabula Muris Senis aging axis compares 3-month and 18-month donors, a larger age span than the RRRM-2 young/old contrast; weak projection could reflect biological nonalignment, age-range mismatch, or both. Sixth, the pathway panel is intentionally small, improving interpretability but limiting coverage.

\subsection{Future directions}

Three analyses would most directly strengthen the manuscript. First, exploratory gene and LIONESS-node decomposition should be added for the secondary resolutions that passed stability. This would identify candidate genes/nodes underlying the ISS-T-aligned remodeling direction without promoting them to primary age-rewiring evidence. Second, OSD-253 should be analyzed in a strain-stratified manner, especially C57BL/6J versus C3H/HeJ, to test whether ECM/fibrosis module activation is attenuated in the TLR4-deficient background. Third, OSD-253 control sensitivity should be presented systematically across original and rerun ground-control choices, with claims separated into recurrence, context, and sensitivity categories.

% =====================================================================
\section{Conclusion}

RRRM-2 is not strong enough, by itself, to establish age-dependent kidney network rewiring at primary module/pathway scale. It is strong enough to provide an arm-specific factorial bridge. When combined with external contrast-level analysis, RRRM-2 shows that the ISS-T flight direction, but not the LAR direction, recurs in OSD-513. The recurring biology is ECM/fibrosis activation with tubular-transport remodeling, not simple aging mimicry. The resulting manuscript is therefore best framed as a multi-resolution, cross-OSDR analysis of ISS-T-aligned kidney remodeling after spaceflight, with WGCNA and LIONESS used as complementary annotation and prioritization layers rather than as overextended standalone evidence.

\section*{Author contributions}
I.S. designed and implemented the analytical framework, conducted the contrast-vector, cross-OSDR, external-aging-axis, LIONESS, WGCNA, and follow-up interpretive analyses, drafted the manuscript, and maintained the source code, tests, and run artifacts. PI authorship pending.

\section*{Competing interests}
The author declares no competing interests.

\section*{Data and code availability}
NASA Open Science Data Repository accessions: OSD-771 (RRRM-2), OSD-513, and OSD-253. Tabula Muris Senis kidney female was used as the external aging reference. Pipeline source, configuration files, run manifests, tests, and manuscript drafts are version-controlled in the project repository. The primary v3 result corresponds to the May 18, 2026 run ending in \code{201823\_2500g}.

% =====================================================================
\begin{thebibliography}{99}\small

\bibitem{beheshti2018microrna} Beheshti A, Ray S, Fogle H, Berrios D, Costes SV. A microRNA signature and TGF-$\beta$1 response were identified as the key master regulators for spaceflight response. \emph{PLOS ONE}. 2018;13(7):e0199621.

\bibitem{benjamini1995controlling} Benjamini Y, Hochberg Y. Controlling the false discovery rate: a practical and powerful approach to multiple testing. \emph{J R Stat Soc Series B}. 1995;57(1):289--300.

\bibitem{chouker2025impact} Chouker A, Gallo A, et al. Impact of spaceflight on endocrine, metabolic and kidney function: current evidence, open issues, and potential countermeasures. \emph{BMC Biology}. 2025;23:71.

\bibitem{cockett1962astronautical} Cockett ATK, Beehler CC, Roberts JE. Astronautical urolithiasis: a potential hazard during prolonged weightlessness in space travel. \emph{Journal of Urology}. 1962;88:542--544.

\bibitem{daSilveira2020comprehensive} da Silveira WA, Fazelinia H, Rosenthal SB, et al. Comprehensive multi-omics analysis reveals mitochondrial stress as a central biological hub for spaceflight impact. \emph{Cell}. 2020;183(5):1185--1201.

\bibitem{finch2025strain} Finch RH, Vitry G, Siew K, Walsh SB, Beheshti A, Hardiman G, da Silveira WA. Spaceflight causes strain-dependent gene expression changes in the kidneys of mice. \emph{npj Microgravity}. 2025;11:11.

\bibitem{galazka2020rnaseq} Galazka JM, et al. RNAseq analysis of rodent spaceflight experiments is confounded by sample collection techniques. \emph{iScience}. 2020;23(12):101733.

\bibitem{kassemi2021numerical} Kassemi M, Brock R, Nemeth N. Numerical characterization of astronaut CaOx renal stone incidence rates to quantify in-flight and post-flight relative risk. \emph{npj Microgravity}. 2021;7:34.

\bibitem{kuijjer2019estimating} Kuijjer ML, Tung MG, Yuan G, Quackenbush J, Glass K. Estimating sample-specific regulatory networks. \emph{iScience}. 2019;14:226--240.

\bibitem{langfelder2008wgcna} Langfelder P, Horvath S. WGCNA: an R package for weighted correlation network analysis. \emph{BMC Bioinformatics}. 2008;9:559.

\bibitem{love2014moderated} Love MI, Huber W, Anders S. Moderated estimation of fold change and dispersion for RNA-seq data with DESeq2. \emph{Genome Biology}. 2014;15:550.

\bibitem{pietrzyk2007renal} Pietrzyk RA, Jones JA, Sams CF, Whitson PA. Renal stone formation among astronauts. \emph{Aviation, Space, and Environmental Medicine}. 2007;78(4 Suppl):A9--A13.

\bibitem{ray2019genelab} Ray S, Gebre S, Fogle H, et al. GeneLab: omics database for spaceflight experiments. \emph{Bioinformatics}. 2019;35(10):1753--1759.

\bibitem{ritchie2015limma} Ritchie ME, Phipson B, Wu D, et al. limma powers differential expression analyses for RNA-sequencing and microarray studies. \emph{Nucleic Acids Research}. 2015;43(7):e47.

\bibitem{robinson2010edger} Robinson MD, McCarthy DJ, Smyth GK. edgeR: a Bioconductor package for differential expression analysis of digital gene expression data. \emph{Bioinformatics}. 2010;26(1):139--140.

\bibitem{siew2024cosmic} Siew K, Nestler KA, Nelson C, et al. Cosmic kidney disease: an integrated pan-omic, physiological and morphological study into spaceflight-induced renal dysfunction. \emph{Nature Communications}. 2024;15:4923.

\bibitem{smith2014men} Smith SM, Heer M, Shackelford LC, Sibonga JD, Spatz J, Pietrzyk RA, Hudson EK, Zwart SR. Men and women in space: bone loss and kidney stone risk after long-duration spaceflight. \emph{Journal of Bone and Mineral Research}. 2014;29(7):1639--1645.

\bibitem{tabula2020} Tabula Muris Consortium. A single-cell transcriptomic atlas characterizes ageing tissues in the mouse. \emph{Nature}. 2020;583(7817):590--595.

\bibitem{tesson2010diffcoex} Tesson BM, Breitling R, Jansen RC. DiffCoEx: a simple and sensitive method to find differentially coexpressed gene modules. \emph{BMC Bioinformatics}. 2010;11:497.

\end{thebibliography}

\vspace{1em}
\begin{center}
\rule{0.4\textwidth}{0.4pt}
\end{center}
\vspace{0.3em}
\noindent\textit{Draft prepared May 2026 for senior-author review. This v3 draft reframes the prior WGCNA-centered manuscript as a multi-resolution cross-OSDR kidney remodeling analysis.}

\end{document}
